% META: {"id": "1NSI-LANG-ME-002", "chapitre": "1NSI-LANGAGE", "type_objet": "methode", "capacites": ["P-LANG-03A", "P-LANG-03B", "P-LANG-03C"], "capacites_codes": ["C3", "C6", "C7"], "methodes": ["M2"], "mode_creation": "ex_nihilo", "sources_inspiration": [], "status": "needs_review"}
% BEGIN-VERIFY
% def moyenne(notes):
%     """Renvoie la moyenne des notes.
%     Precondition : notes est une liste non vide de nombres.
%     Postcondition : le resultat est compris entre min(notes) et max(notes).
%     """
%     return sum(notes) / len(notes)
% r = moyenne([8, 12, 16])
% assert r == 12
% assert min([8, 12, 16]) <= r <= max([8, 12, 16])
% r2 = moyenne([5])
% assert r2 == 5 and min([5]) <= r2 <= max([5])
% try:
%     moyenne([])
%     precondition_violee_detectee = False
% except ZeroDivisionError:
%     precondition_violee_detectee = True
% assert precondition_violee_detectee
% END-VERIFY
\begin{fichemethode}{M2}{Spécifier une fonction : prototype, préconditions, postconditions}

\textbf{Quand l'utiliser ?} Quand l'énoncé demande de « prototyper », de « spécifier », ou
d'énoncer ce qu'une fonction exige et ce qu'elle garantit — avant d'écrire la moindre ligne
de corps.

\textbf{Pas à pas :}
\begin{enumerate}
  \item \textbf{Écris le prototype} : le nom de la fonction, ses paramètres, et ce qu'elle
        renvoie. Le prototype dit \textbf{ce que la fonction offre}, jamais comment elle s'y
        prend. Il précède le corps.
  \item \textbf{Énonce les préconditions} : ce que \textbf{l'appelant doit garantir} avant
        l'appel. Type et nature des arguments, valeurs interdites, tailles minimales.
        Une précondition contraint \emph{l'entrée}.
  \item \textbf{Énonce les postconditions} : ce que \textbf{la fonction garantit} en retour,
        si les préconditions sont respectées. Une postcondition décrit \emph{la sortie}.
  \item \textbf{Vérifie l'orientation de chaque phrase} : « la liste ne doit pas être vide »
        contraint l'appelant, c'est une précondition. « le résultat est compris entre le
        minimum et le maximum » décrit la sortie, c'est une postcondition.
\end{enumerate}

\textbf{Exemple rédigé.}

\begin{python}
def moyenne(notes):
    """Renvoie la moyenne des notes.
    Precondition : notes est une liste non vide de nombres.
    Postcondition : le resultat est compris entre min(notes) et max(notes).
    """
    return sum(notes) / len(notes)
\end{python}

Sur \lstinline{[8, 12, 16]} la fonction renvoie \lstinline{12}, qui est bien compris entre
\lstinline{8} et \lstinline{16} : la postcondition tient. Sur la liste vide, la précondition
est violée et l'exécution échoue — ce n'est pas un défaut de la fonction, c'est un appel
hors contrat.

\textbf{Erreur fréquente.} Décrire l'implémentation au lieu du contrat : « la fonction fait une
boucle sur la liste » ne dit rien à l'appelant. Et la longueur du code n'engage personne : elle
n'apparaît dans aucune spécification.

\end{fichemethode}
